\scanpage{84}
Enfin, si $X=\bigoplus_{1\leq i\leq r}X_i$, la décomposition
de $\Gamma^*(X)$ en \deleted{les} $\bigotimes\Gamma^*(X_i)$ donne
\deleted{une} décomposition des $P_n(\prod X_i,Y)$ en
des $P_{\nu_1\ldots\nu_r}(X_1,\ldots,X_r,Y)$.

On vérifie alors les axiomes
a) à g), et le fait qu'\deleted{il y a}
les produits tensoriels de $P_1$
sont \uline{stricts} \unsure{que|qq} $\unsure{\theta|0|\otimes}$ (nb quelcon. de
facteurs), que les $\Gamma^{\nu_i}$ existent et
satisfont la condition énoncée précédemment \dots.

\medskip\noindent\hfil\rule{4cm}{0.4pt}\hfil\par

Une catégorie \unsure{polynomiale|polynômiale} \inserted{$P$} est dite
\uline{stricte} si elle satisfait aux
conditions préliminaires permettant
de définir $\otimes$ et $\Gamma^*$ avec toutes les
propriétés envisagées. D'autre part,
une cat additive $P_{1}$ avec $\otimes$ et $\Gamma^*$
avec les structures décrites, est
dite une $\otimes$-pd-catégorie. Les
deux notions sont donc en \unsure{équivalence}.

On dit que $P$ est \deleted{\unsure{$k$-li}} une
catégorie polynômiale (resp $\otimes$-pd-cat)
$k$-linéaire ($k$ anneau comm)
si on a donné un hom d'anneau
\[
k\longrightarrow\operatorname{End}(\unsure{\mathrm{Id}|\mathrm{id}}_{P_1})\quad(\subset\mathrm{id}_P)
\]
% archive sheet number: 83
